%% IEEE Conference Paper -- BDHS Maternal and Child Health Risk Prediction
%% Compile with pdfLaTeX. The source follows the supplied IEEE conference template.

\documentclass[conference]{IEEEtran}
\IEEEoverridecommandlockouts
% IEEE requires Times. IEEEtran asks the class for ptm, but under
% XeTeX/Tectonic that shape is undefined in the TU encoding and LaTeX
% silently falls back to Latin Modern. Asking explicitly works on every
% engine, pdfLaTeX included.
\usepackage{newtxtext}
% Splits the reference list across both columns of the last page, which
% otherwise ends after 45% of the left column and leaves the right empty.
\usepackage{balance}
\usepackage{cite}
\usepackage{amsmath,amssymb,amsfonts}
\usepackage{graphicx}
\usepackage{textcomp}
\usepackage{xcolor}
\usepackage{booktabs}
\usepackage{multirow}
\usepackage{url}
\usepackage[hidelinks]{hyperref}
\def\BibTeX{{\rm B\kern-.05em{\sc i\kern-.025em b}\kern-.08em
    T\kern-.1667em\lower.7ex\hbox{E}\kern-.125emX}}

\begin{document}

\title{Prediction of Maternal and Child Health Outcomes from the Bangladesh
Demographic and Health Survey}

\author{\IEEEauthorblockN{MD Farhan}
\IEEEauthorblockA{\textit{Department of Computer Science and Engineering} \\
\textit{Brac University}\\
Dhaka, Bangladesh \\
md.farhan@g.bracu.ac.bd}
\and
\IEEEauthorblockN{SM Sakib}
\IEEEauthorblockA{\textit{Department of Computer Science and Engineering} \\
\textit{Brac University}\\
Dhaka, Bangladesh \\
sm.sakib@g.bracu.ac.bd}}

\maketitle

\begin{abstract}
Bangladesh has made substantial progress against maternal and child mortality
over the past three decades, but the remaining risk is still considerable and
unevenly distributed. Neonatal death occurs in 4\% of live births, 28.7\% of
measured young children are stunted, and only 57.2\% of recent births occur in
a health facility. Those most exposed to these outcomes are often the least
able to reach care promptly. At an initial community-health visit, the worker
usually has no laboratory measurements; the available information is closer
to the household and reproductive history already captured by a national
survey. We examine how useful that information is for early risk screening.
The 2022 and 2017--18 Bangladesh Demographic and Health Surveys are combined
into a leakage-audited dataset of 121,067 records, and five outcomes covering
delivery, the neonatal period and early childhood are predicted. We compare
Logistic Regression, Random Forest, XGBoost, LightGBM, CatBoost, RealMLP and
TabM, together with a greedy post-hoc ensemble and the pretrained transformer
TabPFN-3. All splits are grouped by maternal identity. Test ROC-AUC is between
0.688 for underweight and 0.845 for facility delivery; the methods differ by
no more than 0.07 on any target. Two checks temper this result. A model trained
on 2017--18 preserves much of its ranking performance on 2022 but overpredicts
prevalence by as much as 5.9 times, while sensitivity at a 30\% flag rate
differs 5.7-fold across wealth quintiles. The limiting factor therefore seems
to be the information recorded in the survey rather than the choice of model.
\end{abstract}

\begin{IEEEkeywords}
maternal and child health, Demographic and Health Survey, clinical prediction,
data leakage, temporal validation, algorithmic fairness, tabular learning,
CatBoost, TabM, ensemble selection
\end{IEEEkeywords}

\section{Introduction}

Bangladesh has substantially reduced maternal and under-five mortality, but
the remaining burden is not shared evenly. Across the two survey rounds used
here, neonatal death occurs in 4.0\% of live births, 28.7\% of measured
children are stunted, and 57.2\% of recent births take place in a health
facility. Facility delivery ranges from 34.0\% in the poorest wealth quintile
to 82.6\% in the richest.

During an initial antenatal contact, a community health worker can observe
much the same broad characteristics represented in a household survey: age,
parity, birth spacing, education, wealth, water and sanitation, and previous
reproductive history. Laboratory measurements and biomarkers are not assumed
to be available. The question is whether these variables can provide a useful
early ranking of risk and retain that usefulness across survey rounds and
socio-economic groups.

Many Bangladesh DHS prediction studies use complete cases from a single
module, focus on one outcome, and rely on random splits without formally
checking for leakage. Near-perfect discrimination is hard to reconcile with
the pooled accuracy reported in a recent DHS meta-analysis \cite{rao2025}.
Here we use one harmonised pipeline for two rounds, split records by mother,
select targets according to observed learnability, test transport over time,
and examine performance by wealth stratum.

The study makes four contributions: (i) a pooled corpus checked against
round-specific indicators; (ii) a five-part leakage taxonomy with measurements
before and after each control; (iii) a comparison of seven current tabular
learners plus a post-hoc ensemble; and (iv) operational and equity analyses at
clearly defined screening thresholds.

\section{Related Work}

Eight studies provide the context: six use Bangladesh DHS data directly, one
combines six East African countries and several survey rounds, and one is a
meta-analysis of expected performance. Seven appeared in 2024 or later.

\subsection{Prediction Studies on Bangladesh DHS Data}
Salam \textit{et al.} \cite{salam2025} used ten algorithms to predict four or
more antenatal visits for 5,128 BDHS 2022 women. A random-forest variant led
with 69.5\% accuracy and AUC 0.734, providing the closest external comparison
for our facility-delivery target.

Noor \textit{et al.} \cite{noor2025} predicted the maternal continuum of four
antenatal visits, skilled facility delivery and postnatal contact within 48
hours. In the same 5,128 women, random forest reached 82\% accuracy and AUC
0.889, while only 25.3\% completed the continuum. We instead model stages
separately, allowing the likely point of drop-out to be identified.

Zoha \textit{et al.} \cite{zoha2025} predicted caesarean section from 2,490
BDHS 2022 records and reported AUC 0.871 with 45.6\% prevalence, well above
the 10--15\% considered clinically indicated by the World Health Organization.
Because this signal may reflect health-system behaviour rather than obstetric
need, we fitted but excluded the outcome from the final panel
(Section~\ref{sec:targets}).

Sultana \textit{et al.} \cite{sultana2025} compared six algorithms on 3,192
complete BDHS 2022 mother--child records. XGBoost led at 80\% accuracy and AUC
0.761, with SHAP highlighting gestational age and administrative division. The
study also illustrates the attrition caused by requiring measured birth weight.

Naznin \textit{et al.} \cite{naznin2025} reported 98.75\% accuracy and 99.79\%
AUROC for under-five mortality among 4,913 BDHS 2022 respondents. Such
discrimination is unusual for socio-demographic predictors and can indicate
pre-split resampling or post-outcome features. We therefore treat it as a
cautionary comparison; it also motivated our leakage audit, where neonatal
death falls from ROC-AUC 0.967 to 0.773 after removing an arithmetic identity.

Islam \textit{et al.} \cite{islam2024} modelled stunting, wasting and
underweight in BDHS 2017--18. XGBoost reached 81.7\% accuracy and AUC 0.802
for stunting, but 76.2\% accuracy and AUC 0.622 for wasting. This contrast is
why we use PR-AUC rather than headline accuracy for model selection.

\subsection{Pooling Rounds, and the Performance Envelope}
Tadese \textit{et al.} \cite{tadese2024} pooled 27,691 children from DHS rounds
between 2016 and 2021 across six East African countries to predict incomplete
immunisation. XGBoost led at 79.0\% accuracy and AUC 0.86. The multi-round
design and mostly maternal/reproductive predictors are directly relevant here.

A meta-analysis of eleven DHS malnutrition studies \cite{rao2025} reports pooled
accuracies of 68.92\% for stunting, 84.39\% for wasting and 73.60\% for
underweight, with substantial heterogeneity. These values set a realistic
reference for survey-based prediction and show that a connected
maternal-and-child panel remains relatively unexplored.

\subsection{Tabular Modelling Practice}
Gradient-boosted trees remain strong general-purpose baselines
\cite{xgboost,lightgbm,catboost}. TabArena \cite{tabarena} benchmarks sixteen
models on fifty-one datasets, ranks CatBoost first among individual learners,
and finds that diverse post-hoc ensembles can improve on single models. This
motivates our greedy ensemble (Section~\ref{sec:models}), following Caruana
\textit{et al.} \cite{caruana2004}. RealMLP and TabM are newer neural tabular
models \cite{realmlp,tabm}; TabPFN \cite{tabpfn} is a transformer pretrained
on synthetic tabular tasks. TRIPOD+AI \cite{tripod} and equality-of-opportunity
criteria \cite{hardt2016} motivate our external and subgroup checks.

\subsection{Identified Gap}
The reviewed studies leave five gaps: connected maternal-and-child prediction,
round pooling under one pipeline, leakage auditing, treatment of clustering,
and wealth-stratified performance. We address all five and treat the resulting
lower but defensible performance as a finding rather than a failure.

\section{Dataset and Exploratory Analysis}

\subsection{Source and Nested Cohorts}
We use BDHS 2022 (DHS-8) and BDHS 2017--18 (DHS-7) \cite{bdhs2022}. Their
Individual and Birth recodes were merged through one pipeline, producing
121,067 records, 899 variables, 1,346 clusters and 45,458 mothers. Since DHS
modules are not observed for every record, analyses use nested cohorts: broad
mother variables, narrower birth histories, recent-birth maternity variables,
and anthropometry for physically measured children.

\begin{table}[t]
\caption{Analysis cohorts and targets}
\label{tab:cohorts}
\centering
\small
\resizebox{\columnwidth}{!}{%
\begin{tabular}{@{}llr@{}}
\toprule
\textbf{Cohort} & \textbf{Targets} & \textbf{Records} \\
\midrule
births & neonatal death & 112,546 \\
mother & equity analysis & 45,458 \\
child & stunting, severe stunting, underweight & 11,954 \\
recent birth & facility delivery & 10,773 \\
\bottomrule
\end{tabular}}
\end{table}

\begin{figure*}[t]
\centering
\includegraphics[width=0.78\textwidth]{figs/outcome_prevalence.png}
\caption{Outcome prevalence with 95\% design-based intervals. The broad range
of base rates motivates target-specific metrics and thresholds.}
\label{fig:eda}
\end{figure*}

Figure~\ref{fig:eda} shows outcomes ranging from common pathway measures to
rare adverse events, so accuracy alone is misleading. Missingness is treated
as survey structure within each cohort, not as an automatic reason to remove a
column. Ten published indicators were recomputed by round; all were within
four percentage points of the published values. Stunting, for example, was
23.22\% versus 23.6\% in 2022 and 30.72\% versus 31.0\% in 2017--18.

\subsection{Targets and Selection}
\label{sec:targets}
The final panel contains facility delivery, neonatal death, stunting,
underweight and severe stunting, with positive rates of 57.2\%, 4.0\%, 28.7\%,
22.4\% and 7.9\%. Of fifteen candidates, fever, wasting and diarrhoea were
difficult to learn; caesarean section reflected clinical decision-making; and
pregnancy loss depended on recall timing. Incomplete immunisation fell from
ROC-AUC 0.944 to 0.59 after restricting to the 12--23 month denominator and
removing age, exposing age leakage.

\section{Methodology}

\subsection{Cleaning and Leakage}
\label{sec:leakage}
DHS files combine valid measurements, meaningful zeros and nonresponse codes.
We decode them into missing values, true zeros and common numeric units, and
namespace cluster and maternal identifiers by round. Feature blocks cover
demographic, socio-economic, empowerment, household, reproductive,
healthcare-access and child variables; 96--118 features remain per target.

Five leakage forms are audited: post-outcome variables, same-moment
co-measurements, definitional variables, arithmetic identities, and informative
missingness. For neonatal death, retaining the fertility identity gives
ROC-AUC 0.967, versus 0.773 after blocking it; a survivor-only age variable
gives 0.987. Both values and observation patterns therefore require inspection.

\begin{table}[t]
\caption{Sensitivity to leakage controls for neonatal death}
\label{tab:leakage}
\centering
\small
\resizebox{\columnwidth}{!}{%
\begin{tabular}{@{}lrl@{}}
\toprule
\textbf{Feature set} & \textbf{ROC-AUC} & \textbf{Audit interpretation} \\
\midrule
Leakage-audited & 0.773 & Intended early-risk estimate \\
Fertility identity retained & 0.967 & Arithmetic label proxy \\
Survivor-only age retained & 0.987 & Missingness is the label \\
\bottomrule
\end{tabular}}
\end{table}

\subsection{Feature Representation}
Boosted trees use missingness natively. Linear and neural paths use median
imputation with indicators, robust scaling and rare-level pooling; ordered DHS
scales are ordinal-encoded. Splits are grouped by maternal identity because
90\% of birth rows belong to mothers appearing more than once.

\subsection{Sampling and Outcome Construction}
Each target is constructed only within its valid DHS universe: facility
delivery among births in the preceding five years, neonatal death from birth
history, and anthropometric outcomes among children with DHS height and weight
measurements. An unadministered module is not treated as a negative response.
Feature selection follows the same cohort-specific rule.

Sampling weights estimate national prevalence and design-based intervals;
cluster identifiers are retained for uncertainty checks. Held-out records are
used for ROC-AUC and PR-AUC, so weights do not change groups or reveal outcomes
during training. The average design effect is 2.19, making independent-row
intervals about 1.48 times too narrow. Round-specific reconstruction guards
against a pooled prevalence representing neither survey population.

\subsection{Models and the Ensemble}
\label{sec:models}
Seven learners are benchmarked: logistic regression; random forest; XGBoost;
LightGBM; CatBoost \cite{xgboost,lightgbm,catboost}; RealMLP; and TabM
\cite{realmlp,tabm}. TabPFN-3 \cite{tabpfn} is evaluated separately with the
same split, preprocessing and row cap. All learners use published defaults;
there is no hyperparameter search, so the comparison is between model families.

\textbf{What ``Ensemble'' means here.} This term denotes a defined blend. After
fitting the seven learners, greedy forward selection with replacement
\cite{caruana2004} runs for fifteen rounds, appending the already-fitted learner
that most improves PR-AUC on an inner validation split. A learner selected $k$
times receives weight $k/15$; no new model or test-set weight is fitted.
Table~\ref{tab:blend} reports the target-specific compositions.

Operating thresholds minimise an explicit false-negative cost while respecting
a maximum flag rate of 30\%, used here as a screening-capacity constraint.

\begin{table}[t]
\caption{Composition of the winning ensembles}
\label{tab:blend}
\centering
\small
\resizebox{\columnwidth}{!}{%
\begin{tabular}{@{}lrrrrrr@{}}
\toprule
\textbf{Target} & \textbf{RF} & \textbf{XGB} & \textbf{LGBM} &
\textbf{CatB} & \textbf{RealMLP} & \textbf{TabM} \\
\midrule
facility delivery & 1 & 2 & 1 & 5 & 6 & --- \\
neonatal death    & --- & --- & 8 & --- & 6 & 1 \\
severe stunting   & 4 & --- & 2 & 7 & 2 & --- \\
\bottomrule
\end{tabular}}
\par\vspace{2pt}\parbox{\columnwidth}{\footnotesize Entries are selection counts out of fifteen
greedy rounds, so each row sums to fifteen; a dash means the learner was never
selected and carries weight zero. Stunting and underweight are absent because
a single model won those targets outright.}
\end{table}

\subsection{Validation and Extensions}
Internal evaluation uses grouped train--test splits and 400 bootstrap resamples
for confidence intervals; temporal evaluation trains on 2017--18 and tests on
2022. Conformal prediction uses a held-out calibration split, survival analysis
uses person-period rows, and multi-task learning jointly predicts the three
nutrition outcomes. These are secondary analyses.

Figure~\ref{fig:pipeline} outlines the analysis from survey recodes to
cohort-specific risk scores.

\begin{figure}[t]
\centering
\includegraphics[width=\columnwidth]{figs/fig_pipeline.png}
\caption{Analysis pipeline. Survey recodes are decoded once, then passed
through cohort-specific feature and leakage rules before grouped model fitting
and risk reporting.}
\label{fig:pipeline}
\end{figure}

\section{Results}

\subsection{Internal Performance}
Table~\ref{tab:results} lists the strongest model per target. ROC-AUC ranges
from 0.688 to 0.845, with PR-AUC interpreted against each base rate. The full
matrix in Table~\ref{tab:models} shows close model families: the eight
approaches span 0.025 ROC-AUC for facility delivery, 0.030 for stunting and
0.070 for neonatal death. In the last case, random forest's recall falls to
0.001 at the default cut-point on a 4\% base rate, indicating a threshold issue
rather than failed ranking. The earlier nine-model panel changed no target by
more than 0.02 ROC-AUC.

Two winners differ from the ROC-AUC leaders because selection uses PR-AUC.
RealMLP leads stunting on ROC-AUC by 0.002, while CatBoost leads PR-AUC by
0.005; the ensemble leads underweight on ROC-AUC by 0.003, while TabM leads
PR-AUC by 0.003. At a 22.4\% base rate, PR-AUC is more relevant to screening.

\begin{table}[t]
\caption{Best model per target, with 95\% bootstrap intervals.
\textnormal{\textbf{Ensemble}} is the greedy PR-AUC-weighted blend of the
seven benchmarked learners defined in Section~\ref{sec:models}}
\label{tab:results}
\centering
\small
\begin{tabular}{@{}llll@{}}
\toprule
\textbf{Target} & \textbf{Model} & \textbf{ROC-AUC} & \textbf{PR-AUC} \\
\midrule
facility delivery & Ensemble & 0.845 [0.828, 0.859] & 0.878 \\
neonatal death & Ensemble & 0.773 [0.756, 0.789] & 0.209 \\
stunted & CatBoost & 0.707 [0.687, 0.730] & 0.526 \\
severe stunting & Ensemble & 0.700 [0.660, 0.738] & 0.192 \\
underweight & TabM & 0.688 [0.661, 0.712] & 0.415 \\
\bottomrule
\end{tabular}
\par\vspace{2pt}\parbox{\columnwidth}{\footnotesize Winners are selected on PR-AUC. The
composition of each ensemble is given in Table~\ref{tab:blend}; stunting and
underweight were won by single models.}
\end{table}

\begin{table}[t]
\caption{Test ROC-AUC by learner and target.
\textnormal{\textbf{Ensemble}} is not an eighth learner: it is the greedy
PR-AUC-weighted blend of the seven rows above it}
\label{tab:models}
\centering
\small
\resizebox{\columnwidth}{!}{%
\begin{tabular}{@{}lrrrrr@{}}
\toprule
& \textbf{Facility} & \textbf{Neonatal} & & \textbf{Severe} &
\textbf{Under-} \\
\textbf{Learner} & \textbf{delivery} & \textbf{death} & \textbf{Stunted} &
\textbf{stunting} & \textbf{weight} \\
\midrule
Logistic Reg. & 0.839 & 0.719 & 0.686 & 0.687 & 0.681 \\
Random Forest & 0.822 & 0.703 & 0.683 & 0.679 & 0.680 \\
XGBoost       & 0.820 & 0.742 & 0.678 & 0.671 & 0.636 \\
LightGBM      & 0.840 & 0.771 & 0.703 & 0.665 & 0.672 \\
CatBoost      & 0.843 & 0.748 & 0.706 & 0.698 & 0.679 \\
RealMLP       & 0.841 & 0.750 & \textbf{0.708} & 0.696 & 0.687 \\
TabM          & 0.842 & 0.707 & 0.702 & 0.675 & 0.687 \\
\midrule
\textbf{Ensemble} & \textbf{0.845} & \textbf{0.773} & 0.705 &
\textbf{0.700} & \textbf{0.690} \\
\midrule
range         & 0.025 & 0.070 & 0.030 & 0.035 & 0.054 \\
\bottomrule
\end{tabular}}
\par\vspace{2pt}\parbox{\columnwidth}{\footnotesize The ensemble is defined in
Section~\ref{sec:models} and composed in Table~\ref{tab:blend}. Best per
column in bold; the last row is the spread between the largest and smallest
value above it. Values are point estimates on the held-out split, whereas
Table~\ref{tab:results} reports bootstrap medians, so the two can differ in
the third decimal.}
\end{table}

At cost-sensitive operating points, recall increases from 0.167 to 0.502 for
underweight and from 0.123 to 0.425 for severe stunting, with more false
alarms. Split conformal prediction provides 89.3--90.9\% coverage at nominal
90\%.

\subsection{A Pretrained Transformer Changes Little}
\label{sec:tabpfn}
TabPFN-3 was tested with the same grouped split, native-missingness
preprocessing and 10{,}000-row context cap. Its PR-AUC values were 0.872,
0.166, 0.516, 0.185 and 0.417 for facility delivery, neonatal death,
stunting, severe stunting and underweight. It trails the incumbent on three
targets and leads on two by 0.0005 and 0.0046 PR-AUC, only 0.5\% and 5.6\% of
the corresponding interval widths. Its severe-stunting ROC-AUC is 0.711, but
the PR-AUC ranking is unchanged. The same performance range under a different
inductive bias supports an informational, rather than architectural, ceiling.

\subsection{Robustness and Secondary Analyses}
The ensemble improved three targets by at most 0.002 and was slightly worse on
two. This agrees with Table~\ref{tab:models}: once the data set the error
ceiling, model diversity leaves little error to cancel. A tuned nine-model
panel previously reached ROC-AUC 0.848, 0.769, 0.694, 0.699 and 0.678 for the
five targets; the current untuned panel reached 0.845, 0.773, 0.707, 0.700 and
0.688. The change in algorithm panel does not alter the conclusion.

The person-period formulation expanded 104,025 births into 388,016 rows. The
estimated mortality hazard was 3.99\% in month one, 1.49\% in months 1--12,
0.46\% in months 12--24 and 0.65\% in months 24--60. A child observed at 30
months is censored after three completed intervals, not treated as a five-year
survivor. Multi-task learning improved the no-sharing ablation by 0.010 mean
ROC-AUC but remained below single-task CatBoost.

\subsection{Temporal Validation}
Models trained on 2017--18 preserve much of their ranking ability on 2022, but
their probabilities do not transport well (Table~\ref{tab:temporal}).

\begin{table}[t]
\caption{Temporal validation: train 2017--18, test 2022}
\label{tab:temporal}
\centering
\small
\begin{tabular}{@{}lrrrr@{}}
\toprule
\textbf{Target} & \textbf{Int.} & \textbf{Temp.} & \textbf{$\Delta$} & \textbf{Pred./Obs.} \\
\midrule
facility delivery & 0.839 & 0.796 & $-0.043$ & 0.88 \\
neonatal death & 0.766 & 0.687 & $-0.079$ & 5.90 \\
stunted & 0.708 & 0.701 & $-0.007$ & 1.57 \\
severe stunting & 0.686 & 0.682 & $-0.004$ & 2.81 \\
underweight & 0.677 & 0.672 & $-0.005$ & 1.44 \\
\bottomrule
\end{tabular}
\end{table}

\begin{figure}[t]
\centering
\includegraphics[width=\columnwidth]{figs/fig_temporal.png}
\caption{Calibration drift across survey rounds. Bars compare observed 2022
prevalence with mean risk from a model trained on 2017--18; labels give the
ratios in Table~\ref{tab:temporal}. The companion internal-versus-temporal
ROC-AUC panel was omitted because it repeated columns two and three of the
table.}
\label{fig:temporal}
\end{figure}

The nutrition targets lose only 0.004--0.007 ROC-AUC across rounds, yet
predicted prevalence is 1.44 to 5.90 times observed prevalence. For neonatal
death, mean predicted risk is 21.4\% versus an observed 3.6\%. The model may
still rank cases for triage after transport, but probabilities and thresholds
need recalibration.

\subsection{Subgroup Performance}
\label{sec:fairness}
Figure~\ref{fig:fairness} compares ROC-AUC and recall at a common 30\% flag rate
by wealth quintile. Stunting recall is 0.685 in the poorest and 0.121 in the
richest, a 5.7-fold difference. For facility delivery, recall is 0.119 in the
poorest and 0.812 in the richest. Discrimination is steadier than thresholded
sensitivity, so a global cut-point needs stratum-specific reporting or
threshold stratification \cite{hardt2016}.

\subsection{Operational Implications}
The 30\% flag rate is a capacity scenario, not a clinical recommendation. With
no cap, the tenfold false-negative penalty for neonatal death would flag most
births, which is impractical. Under the cap, the output is a ranked referral
list rather than a diagnosis. A service could recalibrate on a small
current-round sample and choose a threshold with its staffing and referral
capacity.

% Promoted to a full-width float. Two panels side by side inside a single
% column scale the tick and legend text down to about 4 pt -- the same defect
% the pipeline figure had -- and spanning both columns fixes it without
% touching the plot itself.
\begin{figure*}[t]
\centering
\includegraphics[width=0.86\textwidth]{figs/fig_fairness.png}
\caption{Subgroup performance by wealth quintile. Discrimination varies less
than sensitivity at a single global threshold.}
\label{fig:fairness}
\end{figure*}

\section{Discussion}

The main finding is that survey information, not model capacity, limits
performance. The seven linear, tree-based and neural learners remain within
0.07 ROC-AUC on every target. Leakage auditing and maternal-grouped validation
explain more of the gap from implausibly high published values than changing
the learner.

Limitations are structural: DHS lacks blood pressure, glucose and biomarkers;
the records are retrospective; and 65.7\% of the measured-child cohort comes
from 2017--18. Results are associational and need prospective validation.
Future work should add facility supply variables, test recalibration on a
held-out period, and assess whether risk-ranked triage improves outcomes.

Nutrition targets are correlated (pairwise phi 0.39--0.51) and are not
independent clinical syndromes. The survey does not show whether flagged women
received effective interventions, and the transport experiment covers one
historical gap rather than future policy or climate shocks. All learners used
defaults; broader search might improve a score but would not remove the need
for leakage, temporal and subgroup audits.

\subsection{Recommended Reporting Practice}
Future DHS prediction studies should report the valid universe, record count
and positive rate for each target rather than one pooled denominator. Predictors
should be labelled by information time (before contact, same interview, or
after outcome) so leakage is auditable. Repeated mothers, siblings and survey
clusters should also define resampling groups; a random row split is not
independent when households repeat.

Evaluation should include internal and temporal discrimination, calibration or
predicted-to-observed prevalence, and thresholded sensitivity by relevant
socio-economic strata. Global ROC-AUC alone cannot establish safe transport or
equity. Screening reports should state the flag-rate constraint, false-negative
cost, recalibration sample and referral capacity with the threshold.

\subsection{Reproducibility and Audit Trail}
Data preparation, exploration, modelling and testing are separate rerunnable
stages. The cleaned Parquet table, column dictionary, target registry and model
manifest are retained. The manifest records each target's feature block,
leakage blacklist, split groups and model family, so a changed valid universe
changes the cohort definition rather than every model silently.

Before training results are interpreted, EDA compares weighted prevalence with
published round-specific indicators, inspects missingness by module and cohort,
and checks correlations for definitional duplicates. Testing stores
predictions, bootstrap indices, threshold decisions and subgroup labels, so
temporal and wealth analyses reuse the aggregate predictions. These artefacts
also permit one target to be rerun after a new survey round is added.

\section{Conclusion}

Using household survey variables alone, a two-round leakage-audited BDHS
pipeline ranks five maternal and child outcomes with test ROC-AUC from 0.688
to 0.845. Linear, tree-based and neural learners, their greedy ensemble, and a
pretrained tabular transformer all fall within this range, placing the ceiling
in the data rather than the algorithm.

Two audits determine usability. Temporal transport preserves ranking better
than calibration: nutrition targets lose less than 0.007 ROC-AUC, but predicted
prevalence is 1.44--5.90 times observed, requiring current-sample
recalibration. A global threshold also creates a 5.7-fold sensitivity gap
across wealth quintiles, with direction varying by outcome. Thus the practical
output is a ranked referral list with a stated flag rate, false-negative cost
and stratum-specific sensitivity, not a probability taken at face value.

\section*{Acknowledgment}
The authors thank the Demographic and Health Surveys Program and the National
Institute of Population Research and Training, Bangladesh, for providing the
survey data used in this study.

% No \clearpage here. It forced the reference list onto a page of its own and
% left three quarters of the preceding page blank; \balance alone evens the
% two columns of whatever page the list happens to end on.
\balance
\begin{thebibliography}{00}

\bibitem{rao2025} B. Rao \textit{et al.}, ``Machine learning in predicting
child malnutrition: a meta-analysis of Demographic and Health Surveys data,''
\textit{International Journal of Environmental Research and Public Health},
vol. 22, no. 3, p. 449, 2025.

\bibitem{salam2025} M. A. Salam, M. M. Islam, and M. R. Karim, ``Machine
learning algorithms for predicting and identifying the influencing predictors
of antenatal care visits among women in Bangladesh: evidence from BDHS 2022
data,'' \textit{PLOS ONE}, 2025.

\bibitem{noor2025} S. T. A. Noor, R. B. Islam, S. Yeasar, and S. Siddique,
``Machine learning-based prediction of maternal continuum of care completion:
evidence from Bangladesh Demographic and Health Survey 2022,''
\textit{Array}, 2025.

\bibitem{zoha2025} S. Zoha \textit{et al.}, ``Identifying determinants and
predicting cesarean section delivery among Bangladeshi women using machine
learning: insight from BDHS 2022 data,'' \textit{PLOS Global Public Health},
2025.

\bibitem{sultana2025} N. Sultana \textit{et al.}, ``Machine learning based
prediction of low birth weight and its associated risk factors: insights from
the Bangladesh Demographic and Health Survey 2022,'' \textit{PLOS Global
Public Health}, 2025.

\bibitem{naznin2025} S. Naznin \textit{et al.}, ``Identifying determinants of
under-5 mortality in Bangladesh: a machine learning approach with BDHS 2022
data,'' \textit{PLOS ONE}, 2025.

\bibitem{islam2024} M. M. Islam \textit{et al.}, ``Prediction of
undernutrition and identification of its influencing predictors among
under-five children in Bangladesh using explainable machine learning
algorithms,'' \textit{PLOS ONE}, vol. 19, no. 12, e0315393, 2024.

\bibitem{tadese2024} Z. B. Tadese \textit{et al.}, ``Prediction of incomplete
immunization among under-five children in East Africa from recent demographic
and health surveys: a machine learning approach,'' \textit{Scientific
Reports}, vol. 14, 2024.

\bibitem{xgboost} T. Chen and C. Guestrin, ``XGBoost: A scalable tree boosting
system,'' in \textit{Proc. 22nd ACM SIGKDD}, 2016, pp. 785--794.

\bibitem{lightgbm} G. Ke \textit{et al.}, ``LightGBM: A highly efficient
gradient boosting decision tree,'' in \textit{Advances in Neural Information
Processing Systems}, vol. 30, 2017.

\bibitem{catboost} L. Prokhorenkova \textit{et al.}, ``CatBoost: unbiased
boosting with categorical features,'' in \textit{Advances in Neural
Information Processing Systems}, vol. 31, 2018.

\bibitem{tabarena} N. Erickson \textit{et al.}, ``TabArena: a living
benchmark for machine learning on tabular data,'' in \textit{Advances in
Neural Information Processing Systems, Datasets and Benchmarks Track}, 2025.

\bibitem{caruana2004} R. Caruana, A. Niculescu-Mizil, G. Crew, and A. Ksikes,
``Ensemble selection from libraries of models,'' in \textit{Proc. 21st Int.
Conf. Machine Learning (ICML)}, 2004.

\bibitem{realmlp} Y. Gorishniy \textit{et al.}, ``RealMLP: Robust MLP for
tabular data,'' arXiv:2402.06172, 2024.

\bibitem{tabm} Y. Gorishniy \textit{et al.}, ``TabM: Advancing tabular deep
learning with parameter-efficient ensembling,'' in \textit{Proc. ICLR}, 2025.

\bibitem{tabpfn} N. Hollmann \textit{et al.}, ``Accurate predictions on
small data with a tabular foundation model,'' \textit{Nature}, vol. 637,
no. 8045, pp. 319--326, 2025.

\bibitem{tripod} G. S. Collins \textit{et al.}, ``TRIPOD+AI statement:
updated guidance for reporting clinical prediction models that use regression
or machine learning methods,'' \textit{BMJ}, vol. 385, e078378, 2024.

\bibitem{hardt2016} M. Hardt, E. Price, and N. Srebro, ``Equality of
opportunity in supervised learning,'' in \textit{Advances in Neural
Information Processing Systems}, vol. 29, 2016.

\bibitem{bdhs2022} National Institute of Population Research and Training
(NIPORT), Ministry of Health and Family Welfare, Bangladesh, and ICF,
\textit{Bangladesh Demographic and Health Survey 2022}, Dhaka, Bangladesh
and Rockville, MD, USA, 2024.
\end{thebibliography}

\end{document}
